\subsection{Calibration of a terminal marked field on its selected history}
\label{sec:actual-marked-portfolio-calibration}

The local invitation in Subsection~\ref{sec:ex-marked-field-local} is a
terminal-scale construction. The public field calibration also needs one
label correction that works at all earlier scales of the already selected
history. This section follows that exact route: the campaign supplies the
fixed portfolio, the terminal Cream supplies the final coefficient, and a
backward induction through the selected transfer sequence relocates that coefficient.
The history, doors, family, and transfer sequence remain fixed throughout.

\subsubsection{From the campaign to the fixed portfolio}

The campaign produced in Subsection~\ref{sec:marked-campaign-producer} is an
imprint with the powered slot-degree bound. Applying the fixed-stencil
construction of Theorem~\ref{thm:ex-marked-portfolio-section} gives, at
each label degree $d$, one finite section $V_d$, its stencil, its
regrading $R_d$ and its credits. These choices
are made before selecting a finite-volume history. On the public branch the
same data are carried alongside the selected family. They do not replace
the family's own local sections: at degree $d$ the family keeps its fixed
local section $T_d$, which defines its observables $C.O(v)$, and the
portfolio section $V_d$ is compared with it on the selected history. The
comparison is Proposition~\ref{prop:ex-marked-solaris}; its two
ingredients are the one-body invitation of
Proposition~\ref{prop:ex-marked-invitation}, whose hypotheses refer to
the selected family and its terminal Keeper, and the same-history
calibration proved in this section.

The precise same-history obligation is that the credits, budget triple and
terminal Cream of the fixed section $V_d$ are attached to the selected
family and its history. The subsequent argument does not replace
that history by one chosen for each period, degree, or terminal volume.

\subsubsection{Finite backward calibration}

Fix one selected history $H$, period $p$, terminal volume $N$, label frame
$F$, action $A$, section $B$, and a transfer sequence $M$ on this history. At scale $i$,
write $\iota_i:F\to F$ for the transfer sequence index, $\Pi_i(v)$ for its played coefficient,
and $g_i(v)$ for its charge. A state at scale $i$ consists of a primary
$\varpi_i\in F$, a residual coefficient $Z_i$, and a bill $b_i$. Its translated
state at displacement $v$ is
\[
 \mathcal Z_i(v)=Z_i+\Pi_i(\varpi_i+v),\qquad
 c_i(v)=b_i+g_i(\varpi_i+v).
\]
Call the state keyed if $Z_i.x=0$ and $Z_i$ has the prescribed format and
remainder bound. Write $J_i$ for the history's clock index at scale $i$;
the schedule has $J_{i+1}\leq J_i+1$. Assume that the transition from
scale $i+1$ to $i$ preserves keyedness, gives the exact density push
identity for every $v$, and gives the direction relation comparing the
leading coordinates for every $v$. It does not assert equality of those
coordinates.
Writing $L_i$ for the physical side at step $i$,
$e_i=\operatorname{barunit}(L_i)$, and $K_i=W(J_{i+1})$, for each
displacement $v$ there is an invertible map $B_{i,v}$ such that
\[
\begin{gathered}
 \|B_{i,v} x\|_1,\ \|B_{i,v}^{-1}x\|_1\leq K_i\|x\|_1,
 \qquad c_i(v),t_{i+1}(v)\geq0,\\
 \operatorname{core}(B_{i,v}x)=B_{i,v}(\operatorname{core}x),\\
 \|x_i(v)-B_{i,v}(\operatorname{vowel}_{L_i}(x_{i+1}(v)))\|_1
   \leq e_i(\|\operatorname{core}(x_{i+1}(v))\|_1+t_{i+1}(v)),\\
 c_i(v)\leq K_i\|\operatorname{core}(\operatorname{vowel}_{L_i}(x_{i+1}(v)))\|_1
       +e_i t_{i+1}(v).
\end{gathered}
\]
Here $W\geq1$ is a fixed nondecreasing price function on stages,
$\operatorname{barunit}(L)=L^{-3/8}$, $\operatorname{core}x$ is $x$ with
its unit coordinate deleted, and $\operatorname{vowel}_L$ multiplies the
coordinate of a basis vector of grade $(r_j,s_j)$ by $L^{d-r_j}\geq1$.
Further, $x_i(v),x_{i+1}(v)$ are the leading coordinates of the translated
states, and $t_{i+1}(v)=c_{i+1}(v)$ is the upper state's nonnegative bill. More
explicitly, these coordinates are the corresponding transfer sequence-index values
$\iota_i(\varpi_i+v)$ and $\iota_{i+1}(\varpi_{i+1}+v)$; the displayed
direction inequality,
not an equality, compares them. The maps may depend on $v$, but the
displayed constants do not. The density identity is
$\rho_i(v)=L_i^d\,\mathsf P_i\rho_{i+1}(v)$, with $\mathsf P_i$ exactly
the transition push in the history. The finite Tale witness supplies the
weighted support estimate for the specified exponential entries and forced
profiles. Every bound is required at every displacement, not only at the
base displacement $v=0$.

For clarity, the paid one-step input has the following content. At scale
$i$ put $C_i=\operatorname{columns}(i)$. The tuned coupon has a number
$\eta_i\geq0$ with $(1+2C_i)\eta_i\leq e_i$. For every keyed upper
coefficient $U$ with zero leading coordinate, format, and bill $q$, it
provides a finite wing and a coefficient correction whose leading error
cost is at most $\eta_i q$, whose remainder cost is at most $\eta_i q$,
and whose support charge is at most $2q$. The associated Carol witness is
a finite wing that works for every displacement and compares the upper and
lower forced profiles, exponential entries, and charges. The one-step
construction takes the correction coefficient $V$, sets
$h=\iota_i^{-1}V.x$, and uses
\[
 \varpi_i=\varpi_{i+1}+h,\qquad
 Z_i=V+\Pi_i(-h),\qquad
 b_i=\eta_i q+g_i(-h).
\]
The leading coordinate of $Z_i$ vanishes because
$V.x+\iota_i(-h)=0$. The format and remainder bounds follow by adding the
coupon coefficient to the earned played coefficient. For every $v$, the
linear density readout cancels the two terms at $h$ and $-h$; the remaining
density is exactly the transition push of the upper state's density. The
direction estimate follows from the coupon's paid direction bounds, the
transfer sequence cone estimate $g_i(\pm h)\leq C_i\eta_iq$, the Carol charge
comparison, and $(1+2C_i)\eta_i\leq e_i$. Finally, union the two finite
wings and add their Tale inequalities; this gives one common finite support
for all $v$. This is the precise one-step content used in the finite descent.

\begin{lemma}[Backward propagation of one terminal coefficient]
\label{lem:actual-marked-antilles}
Let $0\leq r\leq l\leq m\leq N$. Suppose the history is selected and the
same transfer sequence is admissible and responds at every level from $r$ through $N$.
Suppose its one-step response on each door $D_i$, $l\leq i<m$, is the
response described above: for every keyed input it supplies a keyed output,
the all-displacement density identity and direction bound, and a finite Tale
witness, with the tuning and Coupon bounds stated there. At level $m$, let $Z$ be a
coefficient with the required format, remainder bound $t$, and
\[
             \|Z.x\|_1+t\leq w.
\]
Assume also the scale room and bank budgets. Then there are keyed states
$(\varpi_i,Z_i,b_i)$ for $l\leq i\leq m$ such that:
\begin{enumerate}
\item the terminal state is the state obtained by adjoining $Z$ to the
  level-$m$ transfer sequence and subtracting the transfer sequence's primary which represents
  $Z.x$;
\item every adjacent pair is related by the supplied response on $D_i$;
\item with $\varpi_l$ the bottom primary, evaluate the translated states at
  the displacement $-\varpi_l$: set $x_i=\mathcal Z_i(-\varpi_l).x$ and
  $c_i=c_i(-\varpi_l)=b_i+g_i(\varpi_i-\varpi_l)$. Then, for every
  $l\leq i\leq m$,
  \[
    \exp\!\left(\frac{\sigma_m-\sigma_i}{8}\right)
       (c_i+\|x_i\|_1)\leq 4(1+C_m)w,
    \qquad C_m=\operatorname{columns}(m),
  \]
  where $\sigma_i$ is the history scale parameter;
\item with the same $x_i,c_i$, and with $C_i=\operatorname{columns}(i)$
  and $s_i=L_i^{1/8}$,
  \begin{gather*}
    32(1+C_i)^2\|x_i\|_1\leq c_i\quad(l\leq i\leq m),\qquad
    s_i\|x_{i+1}\|_1\leq c_{i+1}\quad(l\leq i<m),\\
    c_m\leq2(1+C_m)w.
  \end{gather*}
\end{enumerate}
The terminal state has, for every $v\in F$, exact identities
\[
 \mathcal Z_m(v).x=Z.x+\iota_m(v),\qquad
 \mathcal Z_m(v).\mathrm{density}
       =Z.\mathrm{density}+\Pi_m(v).\mathrm{density},
\]
and its terminal bill obeys
\[
 c_m(v)\leq w+C_m\bigl(w+\|\mathcal Z_m(v).x\|_1\bigr),
       \qquad C_m=\operatorname{columns}(m).
\]
\end{lemma}
\begin{proof}
At the terminal level set $a=\iota_m^{-1}(Z.x)$. In the public application
$m=N$, and the top clause of `Sung` gives
$(M.\mathrm{tone}\,N).\mathrm{index}=\mathrm{id}$: its played coefficient
has leading value $v$ at every label $v$. Thus in that application
$\iota_N(v)=v$. Form the residual
coefficient by adding to $Z$ the played field $\Pi_m(-a)$, and set its bill
to $t+g_m(-a)$. Its leading coordinate is zero because
$\Pi_m(-a).x=\iota_m(-a)=-Z.x$. This proves keyedness. For displacement $v$,
linearity of the transfer sequence readout gives
\[
 (Z+\Pi_m(-a)+\Pi_m(a+v)).x=Z.x+\iota_m(v),
\]
and linearity of its density gives cancellation of the $-a$ and $a$
terms, leaving $Z$ plus the density of $\Pi_m(v)$. These are exact
identities.

For the charge bound, the transfer sequence's cone estimate gives
$g_m(x)\leq C_m\|\iota_m(x)\|_1$. Since $\iota_m(a)=Z.x$ and
$\iota_m(a+v)=\mathcal Z_m(v).x$, it follows that
\[
 c_m(v)=t+g_m(-a)+g_m(a+v)
 \leq t+C_m\|Z.x\|_1+C_m\|\mathcal Z_m(v).x\|_1
 \leq w+C_m\bigl(w+\|\mathcal Z_m(v).x\|_1\bigr).
\]
This is the displayed terminal bill bound.

We now descend from $m$ to $l$. At a step $i$, assume a keyed state at
$i+1$ has been constructed. The same transfer sequence supplies its Heard and
Responds data at $i$ and $i+1$. Apply the paid one-step response with the
state's format and remainder bound; its zero leading coordinate is exactly
the Coupon hypothesis. The output is a keyed state at $i$, its transfer sequence-index
formula, the density push identity, the paid direction inequality, and the
finite forced-profile witness. Replace only the state
at $i$ and keep all states at levels above $i$ unchanged. Thus all already
proved chair and transition clauses persist, and the new clauses hold at
$i$. Finite downward induction constructs the full sequence. This is a
finite induction on $m-l$; it does not pass to a subsequence or alter the
selected history.

For the signed difference, let $\varpi_l$ be the primary of the bottom state
and evaluate every state at displacement $-\varpi_l$. The bottom played
coordinate is then zero. At each level, the format and remainder clauses
follow from keyedness and the transfer sequence's earned bounds; transition and density
clauses follow from the all-displacement response identities. For the
numerical propagation fix this displacement $v=-\varpi_l$ and abbreviate
$x_i=x_i(v)$, $c_i=c_i(v)$, and $t_i=t_i(v)$; thus the map used at step $i$
is the possibly displacement-dependent $B_{i,v}$, whose norm bounds are
uniform. Put $C_i=\operatorname{columns}(i)$,
$s_i=L_i^{1/8}$, $e_i=s_i^{-3}$, and $A_i=32(1+C_i)^2$. Room cost says
$64W(J_i+1)^2\leq C_i$; since $J_{i+1}\leq J_i+1$ and $W$ is
nondecreasing, this yields $64K_i^2\leq C_i$. The bank inequality at
step $i$, together with $W(J_{i+1})\leq W(J_{i+1}+1)$, yields
$100000K_i^6(1+C_{i+1})^2\leq s_i$. Also $K_i\geq1$ and
$C_i\geq1$. Consequently
\[
 4K_ie_i\leq1,\qquad 4K_iA_{i+1}e_i\leq1,
       \qquad (4K_i^2+1)(s_ie_i)\leq1.
\]
Indeed, writing $C=C_{i+1}$ and $K=K_i$, the bank bound implies
$s_i\geq128K^2(1+C)^2$, since $100000K^6\geq128K^2$ for $K\geq1$.
Thus $4KA_{i+1}=128K(1+C)^2\leq s_i$ and $4K\leq s_i$.
Also $(4K^2+1)s_ie_i=(4K^2+1)s_i^{-2}\leq1$ because
$s_i\geq128K^2$. These are the three displayed numerical inequalities.
The direction inequalities for the translated states imply, with
$u=\|x_i\|_1$, $u_+=\|x_{i+1}\|_1$, and
$y=\|\operatorname{vowel}_{L_i}(x_{i+1})\|_1$,
\[
 y\leq K_i(u+e_i(y+c_{i+1})),\qquad
 c_i\leq K_i y+e_i c_{i+1},\qquad u_+\leq y.
\]
If $A_i u\leq c_i$, combining these inequalities with the three
numerical bounds gives
$A_{i+1}u_+\leq c_{i+1}$ and $s_i c_i\leq c_{i+1}$.
Indeed, substituting $u\leq(K_i y+e_i c_{i+1})/A_i$ into the first
direction inequality gives
\[
 y\leq \left(\frac{K_i^2}{A_i}+K_ie_i\right)y
       +\left(\frac{K_ie_i}{A_i}+K_ie_i\right)c_{i+1}.
\]
Since $A_i=32(1+C_i)^2$, we have $K_i^2/A_i\leq1/32$;
$K_ie_i\leq1/4$; and $A_i\geq1$. Thus the coefficient of $y$ is at
most $9/32<1/2$, and the coefficient of $c_{i+1}$ is at most
$2K_ie_i$. It follows that $y\leq4K_ie_i c_{i+1}$. Now $u_+\leq y$ and
$4K_iA_{i+1}e_i\leq1$ give $A_{i+1}u_+\leq c_{i+1}$. In the same way
$s_iu_+\leq4K_i(s_ie_i)c_{i+1}\leq c_{i+1}$, because
$4K_i\leq4K_i^2+1$ and $(4K_i^2+1)(s_ie_i)\leq1$; this is the second
inequality of the last item of the lemma. Also
$c_i\leq K_i y+e_i c_{i+1}$ gives
$s_ic_i\leq(4K_i^2+1)s_ie_i c_{i+1}\leq c_{i+1}$ by the third
numerical bound. At the bottom,
keyedness gives $x_l=0$ and nonnegative cost. Thus induction gives
$A_i\|x_i\|_1\leq c_i$ at every level, and iterating
$s_i c_i\leq c_{i+1}$ gives
$\prod_{j=i}^{m-1}s_j c_i\leq c_m$. The scale identity is
$\prod_{j=i}^{m-1}s_j=\exp((\sigma_m-\sigma_i)/8)$.

At the top the terminal bill bound and $A_m\|x_m\|_1\leq c_m$ imply
$2C_m\|x_m\|_1\leq c_m$ and hence
$c_m\leq2(1+C_m)w$. This completes the last item of the lemma.
Since $A_i\geq1$, also $\|x_i\|_1\leq c_i$. Therefore
$\exp((\sigma_m-\sigma_i)/8)(c_i+\|x_i\|_1)
\leq2c_m\leq4(1+C_m)w$. This proves the cone clause directly from
the room and bank inequalities, on the same finite interval and history.
\end{proof}

\subparagraph{Same-history inputs to the lemma.}
For each step, the response is supplied by a same-history Tone record, a
paid Coupon for every keyed silent upper coefficient, and the matching
Carol/Tale witness. The Coupon fixes the input exponential entry, forced
profile, distinguished leading value, and density; its output coefficient
has the specified exact density identity and paid leading/remainder bounds.
The Carol witness includes the upper and lower forced profiles and
exponential entries at every displacement, together with common finite
support and the charge comparison used above. These are exact inputs, not
continuity limits.

The local three-card algebra which produces the paid Coupon from a
Heading/Stampede record is proved in Subsection~\ref{sec:three-port-history}.
It uses the common-history identities, waybill and forced-price estimate
there, together with the three selected card estimates and the finite
threshold maximum. This finite descent consumes the resulting Coupon and
Carol records; it does not reprove the card algebra. The common-history attachment of these records is supplied by
Propositions~\ref{prop:selected-reprise} and~\ref{prop:selected-music},
under the common selected-history purchases. The zero-coordinate
Coupon used in this descent is supplied by
Proposition~\ref{prop:selected-zero-coordinate-coupon}, with its
insertion, soundness and merger hypotheses verified as described in
Subsection~\ref{sec:selected-centered-service}.

\subsubsection{Application to the terminal Cream}

We now state the calibration used by the selected field argument. Fix one
family $C$ which makes the action, its selected history, a period $p$, a
terminal volume $N$, and a degree $d$. Write $a_i$ for the mesh at level
$i$, $L_i=a_i/a_{i+1}$ for the block side, $\sigma_i$ for the running
scale, $C_i=\operatorname{columns}(i)$, $\alpha_i$ for the support
exponent and $B_i$ for the support radius of the level-$i$ budget,
$\theta_i$ for its threshold, $J_i$ for the clock index, and $T$ for the
local section of the family at degree $d$, with frame $F$ and maximal
factor grade $q$. The following notation is used for the root tests:
\begin{align*}
 \operatorname{puffs}(H_i)&=L_i^{1/8},&
 \operatorname{gain}(N,i)&=e^{(\sigma_N-\sigma_i)/8},\\
 \|\operatorname{alto}(H,i,d)\|&=a_i^{-d},&
 \operatorname{escrow}(N,k,E)&=E/\operatorname{gain}(N,k),\\
 \operatorname{stakes}(m,k)&=
   \exp\bigl((a_m/a_k)^4\,\mathrm{Treasury}_k+5\bigr),&
 \operatorname{flight}_{-\gamma}(N,p)&=e^{-\gamma\sigma_N},
\end{align*}
with $\gamma=1/8$ and $\mathrm{Treasury}_k$ the vacuum treasury per site
at level $k$, and
\[
 \operatorname{cider}(d,m,k)=\operatorname{stakes}(m,k)\,
     \frac{4a_k^{-d}}{\zeta_{d,k}}\,e^{\gamma\sigma_k}.
\]
For $m\leq N$, $\operatorname{rootTest}_{m;N,p}(f)$ is the centered root
test of Subsection~\ref{ex-centered-estimates} for the centered tiling of
depth $m$ based at the origin and the finite-volume law $\mu_{N,p}$: the
$K$th root of the modulus of the expectation of the product of the $K$
prescribed copies of $f$, where $K=|\mathrm{Site}(m,p)|$ is the number of
cells. We say that $m$ is \emph{opened at} $k$, and write
$\mathcal O(m,k)$, if $m\leq k$ and $B_k^2\leq a_m/a_k$.

\subparagraph{The paid selected-history record.}
The calibration consumes the following record on the fixed family and
history. The family, the history, the numerical sequences and the indexed
field of doors are fixed globally: they do not depend on $N$, $p$ or $d$,
and the items are required for all of them. The record is assumed in this
section. Its simultaneous production on the selected public run,
including the reserved-address and fee conditions of
Subsection~\ref{sec:terminal-marked-construction}, is supplied under the common selected-history purchases in
Subsection~\ref{sec:selected-centered-service}; its paragraph on the paid
selected-history record identifies each of (R1)--(R7) with its supplier. It consists of a
nondecreasing function $W\geq1$ on stages, real sequences
$u_i,\psi_i,\phi_i,\delta_i$, heights $\zeta_{d,i}$, a sequence of
natural numbers $r_i$, and one gauge door $D_{N,p,i}$ between the
representatives at levels $i+1$ and $i$ for all $i<N$ and all $p$, with
these properties.
\begin{enumerate}
\item[(R1)] \emph{Room and bank.} For every $i$,
  $64W(J_i+1)^2\leq C_i$ and
  $100000\,W(J_{i+1}+1)^6(1+C_{i+1})^2\leq L_i^{1/8}$.
\item[(R2)] \emph{Scale bounds.} For every $i$: $\phi_i>0$,
  $\delta_i>0$, $1\leq\delta_i\phi_ie^{B_i/128}$,
  $B_{i+1}^2\leq L_iB_i^2$ and $8\leq B_i$.
\item[(R3)] \emph{Occupancy.} The sequence $r_i$ is nondecreasing and
  $r_i\geq1$.
\item[(R4)] \emph{Transition service.} Put
  $\mathsf s(B)=e\,(2(16B+9))^4$, and let $w^{\rm f}_i$ be the forced
  allowance at level $i$. For all $N$, $p$ and $i<N$ there are reals
  $y,w$ such that:
  \begin{itemize}
  \item the door $D_{N,p,i}$ satisfies the multi-marker offer input
    with occupancy $r_i$, offer
    threshold $\alpha_{i+1}-4$ and offer allowance $y$, ordinary
    exponent $\alpha_i+\frac12$ with site budget $1$, and flagged
    exponent $\alpha_i+1$ with site budget $w$. Its full conditions
    remain an input; they are not derived here. They include the pushed representation of every finite census of
    offers of occupancy at most $r_i$, soundness and tuning of the door,
    and a contact condition on the upper representation;
  \item $0<\delta_i\leq1$, $\alpha_i\geq100$, $\alpha_{i+1}\geq100$;
  \item with $T=r_i\mathsf s(B_{i+1})\delta_i+4\delta_i$, one has
    $T+e^TT\leq y$;
  \item with $c=w+r_i\mathsf s(B_i)\delta_i$, one has
    $(2B_i+1)^4c\leq\frac12$ and $e^{c}-1\leq1$.
  \end{itemize}
  Moreover, at every level $i$: $u_i\geq0$, $\alpha_i\geq3$, and with
  $C=w^{\rm f}_i+r_i\mathsf s(B_i)u_i$ one has $(2B_i+1)^4C\leq3$ and
  $e^{C}-1\leq1$.
\item[(R5)] \emph{Heights.} For every $i\geq\operatorname{ref}(d)$:
  $\zeta_{d,i}>0$, $\phi_i\leq1$, $\psi_i\geq0$,
  \begin{gather*}
   \phi_iL_i^{d+1}\zeta_{d,i}\leq\zeta_{d,i+1},\qquad
   \phi_iL_i^{d}C_{i+1}\zeta_{d,i}\leq\zeta_{d,i+1},\\
   \zeta_{d,i}C_i\leq1,\qquad
   2\mathcal O_d(\psi_i+2C_{i+1})\zeta_{d,i+1}\leq1,
  \end{gather*}
  where $\mathcal O_d$ is the orbit cost of the frame. For every
  terminal level $N\geq\operatorname{ref}(d)$ and period $p$ there is
  $X\geq1$ with $|T(e_j)(0,U)|\leq X$ for all axes and configurations,
  the footprint of $T$ at the origin has cover number at most one, and
  $e^{\alpha_N+2}X\theta_N^{-q}\zeta_{d,N}\leq u_N$.
\item[(R6)] \emph{The transfer sequence.} For every $N\geq\operatorname{ref}(d)$
  and $p$ there is a transfer sequence $M$ on this history at $(N,p)$ whose played
  coefficients are the entries of the family's transfer data, which is
  admissible at every level from $\operatorname{ref}(d)$ to $N$ and responds
  across every intervening door. Its terminal coefficient family has the properties
  of Proposition~\ref{prop:terminal-marked-voice}: lead $v$ at every
  label $v$, so that the terminal index is the identity, zero regular
  part, charge at most $\|\operatorname{core}v\|_1$, the density identity
  \eqref{eq:terminal-marked-construction-display-T5} of Subsection~\ref{sec:terminal-marked-construction}, and forced
  profile equal to the hook of the singleton response $f_v$. It has a Carol
witness with price $\psi_i$ at every door
  $D_{N,p,i}$, $\operatorname{ref}(d)\leq i<N$: a finite set of wings
  and, for every label $v$, a Tale across the door relating the
  exponential entry and forced profile of the upper played coefficient
  to the forced profile of the lower one, with upper price
  $\mathcal O_d\|\iota_{i+1}(v)\|_1(\psi_i+2C_{i+1})$ and lower price
  $g_i(v)$. The full conditions on the wing and Tale witnesses remain
  inputs here. The transfer sequence may depend on
  $(d,N,p)$.
\item[(R7)] \emph{The paid one-step input.} At every door $D_{N,p,i}$,
  $\operatorname{ref}(d)\leq i<N$, there is a number $\eta_i\geq0$ with
  $(1+2C_i)\eta_i\leq\operatorname{barunit}(L_i)$, the door is sound,
  $\psi_i\geq0$, and there is a finite set of wings such that every
  upper coefficient with zero leading coordinate, the terminal format
  and remainder bill $q$ has a Coupon across the door with budgets
  $(\eta_iq,\eta_iq,2q)$, as described before
  Lemma~\ref{lem:actual-marked-antilles}. The full conditions on the
  Coupon beyond that description remain an input here.
\end{enumerate}
The multi-marker offer input in (R4) has the following interface. Fix
one of its doors, its upper and lower representations, and the associated
parent map from upper to lower sites. Take any finite marker set $V$ and
any downward-closed family of subsets containing the empty set and all
singletons; its marker algebra has commuting square-zero generators, with
only the allowed subsets as modes. Choose sites $o_v$ on the upper lattice,
strictly positive weights $t_v$, and any finite family of offers. Each
offer consists of an upper activity leaf with support $P$, a nonempty
allowed marker mode $A$, and a complex coefficient $c$; its marker sites
$\{o_v:v\in A\}$ lie in $P$. Its fare is the leaf weight times
$|c|\prod_{v\in A}t_v$. The offer is priced if, for every upper site $x$,
the sum of fares over supports containing $x$, each multiplied by
$e^{(\alpha_{i+1}-4)\operatorname{coverNumber}(B_{i+1},P)}$, is at
most $y$. Assume also that at most $r_i$ markers have any one parent
site.

For every such census, weights, sites and priced offer, the input supplies
lower-lattice polynomial families $F_P$ indexed by supports and $K_{P,I}$
indexed by support--inventory pairs. Both are gauge invariant, measurable
and local to their supports. On the lower lattice the marker positions are
the parent sites $b_v=\operatorname{parent}(o_v)$, and each coefficient of a
mode $A$ in these families vanishes unless $\{b_v:v\in A\}$ is contained in
its support. An $F_P$ term vanishes unless $P$ is healthy and has zero
empty-mode coefficient. A $K_{P,I}$ term vanishes unless $I$ is nonempty,
$I\subseteq P$, and the flags in $P$ are exactly $I$. Each family has
nonnegative summable support prices. For each support or support--inventory
pair, there is one nonnegative majorant for each mode that bounds the
absolute value of that mode's coefficient for every configuration. The
weighted sum of these configuration-independent majorants is at most the
corresponding support price; the weight of a mode $A$ is
$\prod_{v\in A}t_v$. At each lower site the
support prices, multiplied by the exponential of the support cover
number, sum to at most $1$ for $F$ with exponent $\alpha_i+1/2$, and at
most $w$ for $K$ with exponent $\alpha_i+1$.

Their assembled lower representation equals exactly the transition push
of the offered upper polynomial. For every lower support--inventory pair
and configuration, the empty-mode coefficient of $K$ is the ordinary
lower activity profile; its singleton coefficient for marker $v$ is the
sum, over the offered leaves whose mode is $\{v\}$, of their coefficients
times their door-transported profiles. These identities hold pointwise,
not merely after integration. The same door is also required to have the
soundness and tuning interfaces: a uniform forced-cost bound and motion
equivariance for its leaf transports, and compatibility of those
transports with every countable grouped leaf sum. Finally, the upper
representation must have the contact addresses: every nonempty support
has a zero-activity regular address, and every nonempty inventory contained
in a support has an address eligible whenever the flags realize that
inventory. These are hypotheses of the present calibration statement. On the
retained public history they are supplied by
Theorem~\ref{thm:selected-retention-public}, with the common purchases
and local analytic hypotheses stated there.

The items (R1)--(R3), (R5) and the inequalities of (R4) are numerical.
The remaining parts of (R4), (R6) and (R7) are the same-history records;
they must hold for the one family and the one indexed field of doors,
for all $d$, $N\geq\operatorname{ref}(d)$ and $p$. For each $(d,N,p)$
one transfer sequence realizing the family's transfer entries is chosen, and that
choice is retained throughout all subsequent scale, census and motion
tests at those indices. Choosing a different history or door field for
each period or terminal volume is insufficient.

\subparagraph{The terminal Cream.}
A terminal Cream at $(N,p,d)$ with field $b$, charge $Q$, bill $t$ and
budget $E$ is a coefficient $Z$ at level $N$ with the fields of
Proposition~\ref{prop:ex-marked-invitation}: the terminal format; the
remainder bound with weight $\sqrt{w_N}$ and bill $t$;
$\|Z.x\|_1+t\leq Q$; $0<Q$; $4(1+C_N)Q\leq E$; an attainment with reserved
mass $B_{\rm cream}$ such that $B_{\rm cream}\zeta_{d,N}\leq u_NE$,
together with the forced-profile hook identity; the density equation
\[
 Z.\mathrm{density}(U)=a_N^d R_N.\mathrm{eval}(U)b(U);
\]
and continuity of $b$. Note that $E>0$.

\begin{proposition}[Same-history field calibration]
\label{prop:actual-marked-figment}
Assume $C.\mathrm{ref}(d)<N$, that the history has the opening schedule
of Subsection~\ref{ex-centered-site-geometry}, and that the paid
selected-history record (R1)--(R7) holds. Let $Z$ be a terminal Cream at
$(N,p,d)$ with field $b$, charge $Q$ and budget $E$. Then there is one
label $v$ such that
\[
 \operatorname{puffs}(H_{N-1})
       \|Z.x-X_N^{-1}v\|_1
       \leq 2(1+C_N)Q,
 \qquad C_N=\operatorname{columns}(N),
\]
and, for every $k$ with $\operatorname{ref}(d)\leq k<N$, every $m$
opened at $k$, and every motion $s$,
\[
 \operatorname{rootTest}_{m;N,p}\bigl(s\cdot(b-C.O(v))\bigr)
       \leq \operatorname{cider}(d,m,k)\,E\,
               \operatorname{flight}_{-\gamma}(N,p).
\]
Here $X_N$ is the
terminal coefficient calibration isomorphism on the selected family, and
$C.O(v)=a_N^{-d}T(\theta_N^{-1}\mathbin{\cdot_F}
X_N^{-1}v)(0,\cdot)$ is its local observable at the origin.
The label $v$ and all bounds use the same $C$, history, and transfer sequence.
\end{proposition}
\begin{proof}
The record (R6) supplies the transfer sequence $M$ and its Heard and Responds fields
from the bottom reference $r=\operatorname{ref}(d)$. Its Carol witnesses,
the Coupons (R7), the budgets (R1), and the terminal Cream bounds are the
inputs of Lemma~\ref{lem:actual-marked-antilles} with $l=r$, $m=N$ and
terminal cone parameter $w=Q$, as specified by the Cream charge bound.
The larger budget $E$ is reserved for the root-test error estimate. This
produces one backward sequence of keyed states $(\varpi_i,Z_i,b_i)$,
$r\leq i\leq N$, on the original history. Let $a=\varpi_r$ be the primary of
its state at $r$, and set $v=X_N(a)$. By invertibility, $X_N^{-1}v=a$.
In the notation of the lemma put $x_i=\mathcal Z_i(-a).x$ and
$c_i=c_i(-a)$.

The terminal coordinate identity of the lemma and $\iota_N=\mathrm{id}$ give
$x_N=Z.x+\iota_N(-a)=Z.x-X_N^{-1}v$. The last item of the lemma gives
$c_N\leq2(1+C_N)Q$ and, at the step $i=N-1$,
\[
 \operatorname{puffs}(H_{N-1})\|x_N\|_1\leq c_N.
\]
Combining these inequalities proves the first displayed estimate. The
step $N-1$ lies in the range of the lemma because
$\operatorname{ref}(d)\leq N-1$, by the explicit assumption
$\operatorname{ref}(d)<N$.

For the root-test estimate fix $k$ and $m$ as in the statement, abbreviate
$\zeta_i=\zeta_{d,i}$, and choose
\[
 w=\frac{\zeta_k/4}
          {\|\operatorname{alto}(H,k,d)\|\,\operatorname{escrow}(N,k,E)}.
\]
Every factor in the denominator is positive. The height record gives
$\zeta_k>0$, the physical scale gives
$\|\operatorname{alto}(H,k,d)\|=a_k^{-d}>0$, and $E>0$ gives positive
escrow. Marched weights are $w_i=w\prod_{q=k}^{i-1}\phi_q$. The exact
recurrence for the physical scale and the gain is
\[
 \|\operatorname{alto}(H,i+1,d)\|
       =L_i^d\|\operatorname{alto}(H,i,d)\|,
 \qquad
 \operatorname{gain}(N,i+1)
       =\operatorname{gain}(N,i)/\operatorname{puffs}(H_i).
\]
Since $1\leq\operatorname{puffs}(H_i)\leq L_i$, the adjacent paid bound
implies inductively
\[
 w_i\|\operatorname{alto}(H,i,d)\|
       \frac{E}{\operatorname{gain}(N,i)}\leq \frac{\zeta_i}{4}
       \qquad (i\geq k).
\]
The base case is equality at $k$ by the definition of $w$. For the step,
the left side is multiplied by at most $\phi_iL_i^{d+1}$; the first
inequality of (R5) pays this by $\zeta_{i+1}/\zeta_i$. Thus the exact
march/height inequality required for the finite root estimate holds at
every later scale.

For the finite-copy construction, at level $i$ the residual coefficient
$Z_i$ of the state is keyed; the coefficient used in the construction is the
translated coefficient family $\mathcal Z_i(-a)=Z_i+\Pi_i(\varpi_i-a)$ with its bill
$c_i=c_i(-a)$. Keyedness of $Z_i$ and the earned bounds of the transfer sequence
transfer format and the remainder bound, with bill $c_i$, to this
translated coefficient family, which need not have zero leading coordinate. The third
item of Lemma~\ref{lem:actual-marked-antilles} with $w=Q$, and the Cream
inequality $4(1+C_N)Q\leq E$, give
$\mathrm{gain}(N,i)(c_i+\|x_i\|_1)\leq E$. Combined with the
per-level price inequality for
\[
 \frac{w_i\|\operatorname{alto}(H,i,d)\|E}{\mathrm{gain}(N,i)},
\]
this gives the finite-copy rate bound
$w_i\|\operatorname{alto}(H,i,d)\|(c_i+\|x_i\|_1)\leq\zeta_i$ by
multiplication. The
Coupon's density debt is exactly the density transition in the finite-copy transition identity;
its Carol witness supplies the accompanying finite Tale support and the
forced-profile/exponential-entry comparison. At the terminal level let
$c=(M.\mathrm{tone}\,N).\mathrm{index}(-a)$. The top clause of `Sung` gives
$(M.\mathrm{tone}\,N).\mathrm{index}=\mathrm{id}$, so $c=-a$; the
terminal spread bound gives $\|c\|_1\leq2E$. Put
$j=\|\operatorname{alto}(H,N,d)\|$ and let $w_N$ be the terminal price.
The height estimate $w_NjE\leq\zeta_N/4$ implies
\[
 2w_Nj\|c\|_1\leq\zeta_N.
\]
The terminal `erected` estimate, applied at this response cost, supplies an
added price $A_{\rm ins}$ with $w_NjA_{\rm ins}\leq u_N/2$. If
$B_{\rm cream}$ is the Cream's original attainment price, the Bank
hypothesis
\[ B_{\rm cream}\zeta_N\leq u_NE, \]
together with
$w_NjE\leq\zeta_N/4$, gives
$w_NjB_{\rm cream}\leq u_N/4\leq u_N/2$. Hence the combined price
$A_{\rm ins}+B_{\rm cream}$ fits the finite-copy allowance $u_N$. The forced profile is
preserved by the Cream hook identity and linearity of
the played forced-profile map. Finally, the finite-copy density identity is the Cream density
equation after the terminal `originate' identity: the two recentering
terms have density equal to the transfer sequence readout, leaving precisely the
terminal marked density. These are all same-history records; no new
terminal volume or history is selected in assembling these data.

The state residuals $Z_i$ are keyed. Their translated coefficient families are the
finite-copy coefficients; the translated coefficient family at the bottom reference has zero
leading coordinate, by the choice of the negative bottom primary.
The Accord identities give the Wilson push and all-displacement profile
updates; the Cream gives the
terminal forced hook and density identity; and the Bank and footprint
cover pay the forced correction. The output is the exact finite-copy identity
for the terminal field $b+\operatorname{initial}(-a)$, the sum of $b$ and
the original-history insertion with primary $-a$. Since the terminal
calibration is linear,
\[
 \operatorname{initial}(-a)= -C.O(X_Na)=-C.O(v),
\]
so that field is $b-C.O(v)$. To spell out the finite root estimate,
fix $m$ opened at $k$ and write $K=|\mathrm{Site}(m,p)|$ for the number
of slots of its census. The finite census estimate needs a site
assignment for these slots at every level from $k$ to $N$. It is the
opened-copy assignment of Lemma~\ref{lem:ex-centered-opened-copy}: the
slots are embedded by the canonical aligned embeddings, they migrate
along the scale embeddings, they are injective, hence of occupancy one,
which is at most $r_k$ by (R3), and distinct slots have separation at
least $B_j^2$ at every level $j$ between $k$ and $N$; at the terminal
level these sites are the centers of the copies in the root test. The
remaining inputs of the finite census estimate are the monotone
occupancy schedule (R3), the scale bounds (R2), the transition service
(R4), the height inequalities (R5), and continuity of the field
$b-C.O(v)$, which holds because $b$ is continuous by the Cream and
$C.O(v)$ is a polynomial cylinder.

\subparagraph{Serving the residual polynomial.}
\label{par:figment-served-residual}
We now construct the representation needed for this census. This step
uses the full transition and profile records in (R4), (R6), (R7), not
only their numerical inequalities. Let $I$ be its finite set of slots,
and use the algebra with commuting generators $z_\nu$, $\nu\in I$,
satisfying $z_\nu^2=0$. Write $F_\nu$ for the prescribed oriented copy of
$b-C.O(v)$. All subsets of $I$ are retained. Define the terminal polynomial
\begin{equation}
 A_N(U)=R_N.\mathrm{eval}(U)\prod_{\nu\in I}(1+z_\nu F_\nu(U)).
 \label{eq:actual-marked-portfolio-calibration-display-star}
\end{equation}
In particular its coefficient at any subset $J\subseteq I$ is
$R_N.\mathrm{eval}\prod_{\nu\in J}F_\nu$. We construct local arrays for
this polynomial; identifying its singleton coefficients alone would not
suffice.

At level $i$ orient the translated coefficient family $\mathcal Z_i(-a)$ at the
site of slot $\nu$ and multiply it by $a_i^{-d}$. Let $g_{i,\nu}$ be
its ordinary exponential entry and $q_{i,\nu}$ its forced profile.
Denote its full density by $\chi_{i,\nu}$. The density decomposition
\eqref{eq:ope-single-coin} expresses $\chi_{i,\nu}$ as
$R_i.\mathrm{eval}\,g_{i,\nu}$ plus the gas with exactly one insertion
of $q_{i,\nu}$. The format and remainder bounds give summable majorants
for these profiles. Linearity, the adjacent density identity, and
$a_{i+1}^{-d}=L_i^d a_i^{-d}$ give
\[
 \chi_{i,\nu}=\mathsf P_i\chi_{i+1,\nu},\qquad
 \chi_{N,\nu}=R_N.\mathrm{eval}\,F_\nu.
\]
The site migration and the soundness of the same door justify these
identities after orientation as well. A conjugating orientation conjugates
all coefficients and responses and preserves their absolute prices.

The sum of the Cream attainment and the negative played attainment
constructed above, rescaled and oriented at slot $\nu$, is a response
$\Psi_\nu(X,U)$ with hook $q_{N,\nu}$. It is local, gauge invariant,
measurable, anchored at that slot, and zero when $X$ is healthy. Its
weighted reserved mass is at most $u_N$. Set
$h_\nu(U)=\sum_X\Psi_\nu(X,U)$. The hook identity, or equivalently the
singleton coefficient of \eqref{eq:transfer-closure}, says that the gas
with one insertion of $q_{N,\nu}$ is
$R_N.\mathrm{eval}\,h_\nu$. Consequently
\[
 R_N.\mathrm{eval}\,F_\nu
   =R_N.\mathrm{eval}\,(g_{N,\nu}+h_\nu).
\]
Telescoping a finite product proves this identity for every product of
slots with the same factor $R_N.\mathrm{eval}$. Thus it holds without
dividing by that factor, even at its zeros. Since the generators have
square zero, \eqref{eq:actual-marked-portfolio-calibration-display-star} is exactly
\begin{equation}
 e^{E_N}\exp\!\left(\sum_\nu z_\nu g_{N,\nu}\right)
 \mathcal G(k_N^0)
 \exp\!\left(\sum_{\nu,X}z_\nu\Psi_\nu(X,\cdot)\right),
 \label{eq:actual-marked-portfolio-calibration-display-double-star}
\end{equation}
where $R_N.\mathrm{eval}=e^{E_N}\mathcal G(k_N^0)$ denotes the actual
background gas, with its background ordinary contribution included in
$E_N$.

Apply the connected-component construction of
\eqref{eq:transfer-closure} to \eqref{eq:actual-marked-portfolio-calibration-display-double-star}. Its vertices are the selected
background forced factors and the response factors $(\nu,X)$; edges
mean intersecting supports. Every nonzero response factor meets a flag,
which is covered by a background forced factor, so every component has
such a background vertex. Replace each component by the product of its
factors, on their union support with their union inventory, and sum
components with that output pair and marker mode. Conversely, components
with disjoint union supports recover uniquely the original selected
factors. This bijection preserves products and all marker modes. The
empty-mode array is $k_N^0$, and the singleton array is the hook
$q_{N,\nu}$: a component with one response consists of that response
and the background factors meeting its support. Remove these prescribed
modes to define the higher flagged array. The ordinary higher array is
zero. In particular we have constructed the higher coefficients rather
than replaced them by products of singleton densities.

For convergence and payment, the anchored estimate
\eqref{eq:transfer-root-hit} converts the per-slot weighted response
mass $u_N$ into site price $r_N\mathsf s(B_N)u_N$ at exponent
$\alpha_N$: the reserved mass is measured at exponent $\alpha_N+2$,
so it pays the one extra exponent in that anchored estimate. The
background forced price $w_N^{\rm f}$ is measured at $\alpha_N$. With
$D_N=w_N^{\rm f}+r_N\mathsf s(B_N)u_N$, (R4) gives
\[
 (2B_N+1)^4D_N\le3,\qquad e^{D_N}-1\le1.
\]
The connected-support bound \eqref{eq:ope-forest-price}, with spare
exponent three, therefore bounds the flagged array at exponent
$\alpha_N-3$ by one. Removing modes cannot increase a price. The same
majorants justify the component regrouping for countable background
addresses by absolute convergence. This proves terminal servedness with
reserve one at pitch $\pi_N=\alpha_N-3$.

Here are the adjacent profile inputs to the descent, explicitly in the
notation of Subsection~\ref{sec:local-transfer}. For each slot at a door
there is a constant $\ell_\nu$, an anchored family of elementary
insertions $M_{\nu j}$ using finitely many slots, and a lower response
$\Psi'_{\nu}$, such that
\[
 \begin{aligned}
 g_{i+1,\nu}-\ell_\nu&=\sum_j z_{\nu j}\operatorname{ord}(M_{\nu j}),&
 q_{i+1,\nu}&=\sum_j z_{\nu j}\operatorname{flag}(M_{\nu j}),\\
 q_{i,\nu}&=\sum_j z_{\nu j}\mathcal Q_i(M_{\nu j})
                         +\mathcal H_{R_i}(\Psi'_\nu).
 \end{aligned}
\]
Here $\mathcal Q_i$ is the same door's quoted profile map and
$\mathcal H_{R_i}$ its lower hook. These are the finite Tale identities
in (R6)--(R7). To obtain them for the translated coefficient family, take the union
of the Coupon wing for the keyed upper residual and the Carol wing
for the played displacement $\varpi_{i+1}-a$, and add the two expansions.
In the lower state the recentering terms cancel:
\[
 Z_i+\Pi_i(\varpi_i-a)
   =V+\Pi_i(-h)+\Pi_i(\varpi_{i+1}+h-a)
   =V+\Pi_i(\varpi_{i+1}-a),
\]
where $V,h$ are the one-step correction and primary defined before
Lemma~\ref{lem:actual-marked-antilles}. The cancellation holds for the
ordinary entry, density, and forced profile separately. Thus the lower
profile in the displayed Tale is exactly the required one. Taking a
union of wings does not change their door, supports or parent map;
the absolute prices add. Before physical rescaling their upper price
is at most
\[
 \mathcal O_d\|x_{i+1}\|_1(\psi_i+2C_{i+1})+2c_{i+1}
 \le\mathcal O_d(\psi_i+2C_{i+1})(c_{i+1}+\|x_{i+1}\|_1),
\]
since $\mathcal O_d,C_{i+1}\ge1$ and $\psi_i\ge0$. More precisely the
Coupon term costs $2b_{i+1}\le2c_{i+1}$. The lower response price at
exponent $\alpha_i+2$ is bounded by
\[
 \eta_i b_{i+1}+g_i(\varpi_{i+1}-a)
 \le\eta_i b_{i+1}+g_i(-h)+g_i(\varpi_{i+1}+h-a)=c_i,
\]
by subadditivity of the played charge and the defining formula for $b_i$.
After rescaling, call
these prices $s_{i,\nu}$ and $t_{i,\nu}$. The required two payments are
\[
 \begin{split}
 2w_{i+1}s_{i,\nu}
 &\le2w_{i+1}a_{i+1}^{-d}\mathcal O_d(\psi_i+2C_{i+1})
             (c_{i+1}+\|x_{i+1}\|_1)\le1,\\
 w_{i+1}t_{i,\nu}&\le w_{i+1}a_i^{-d}c_i
                  \le\zeta_i C_i\le1.
 \end{split}
\]
These follow from the rate bounds, (R5), $w_{i+1}\le w_i$, and
$C_i\ge1$. Both profiles and the upper exponential entry are therefore
paid; a bound on only one profile would not establish this transition.

Now define $A_i=\mathsf P_i A_{i+1}$, acting coefficientwise, and
$\pi_i=\alpha_i$ for $i<N$. Suppose $A_{i+1}$ has the constructed
representation at pitch $\pi_{i+1}$ and reserve one. In particular
$\pi_{i+1}\ge\alpha_{i+1}-4$, including the terminal step. Apply
Subsection~\ref{sec:local-transfer} with upper weights $w_{i+1}$ and
shrinkage $\delta_i$. The Tale just constructed gives singleton cost
$\delta_i$ after real/imaginary compression, while the higher arrays
cost at most $4\delta_i$. Occupancy gives total prepared site cost
$T=r_i\mathsf s(B_{i+1})\delta_i+4\delta_i$.
The compatible-product expansion costs at most $T+e^TT\le y$.
Hence (R4) applies to this entire prepared polynomial, including every
higher mode. Its ordinary and flagged output bounds and
$(2B_i+1)^4c\le\tfrac12$, $e^c-1\le1$ permit the lower response
closure at exponent $\alpha_i+\tfrac12$. The exact pushed singleton
density and the lower hook identify the prescribed $g_{i,\nu}$ and
$q_{i,\nu}$ by \eqref{eq:ope-single-coin}, using the nonzero selected
Wilson density at this scale, whose strict positivity is proved in
Proposition~\ref{prop:selected-reprise}.

Finally, the separation estimate in
Subsection~\ref{ex-centered-construction} gives, for each supported mode
$J$ with $|J|\ge2$,
\[
 \prod_{\nu\in J}w_i
 \le\prod_{\nu\in J}(\delta_iw_{i+1})
       e^{\operatorname{coverNumber}(B_i,X)/2}.
\]
Indeed $w_i\le\delta_iw_{i+1}e^{B_i/128}$ by (R2), and the tour
bound is $B_i|J|\le64\operatorname{coverNumber}(B_i,X)$.
This pays the change back to $w_i$ and leaves exponent $\alpha_i$ and
reserve one. It completes the induction down to $k$. Higher modes
throughout are those of $A_i$, never reset to products of the
single-insertion pushes.

The terminal coefficients are integrable, since the fields are
continuous on the finite-volume compact configuration space. The
bounded measurable transition kernels preserve integrability; Fubini
and their mass identity preserve every coefficient integral and the
common partition normalization. Therefore, for every $J\subseteq I$,
\[
 \frac{\int[z_J]A_k}{\int R_k.\mathrm{eval}}
       =\int\prod_{\nu\in J}F_\nu\,d\mu_{N,p}.
\]
All integrals on the left use the respective normalized priors. In
particular the full mode is the original residual census moment.
We have proved the required served representation at $k<N$ with
exponent $\alpha_k$ and reserve one. With constant positive slot weight
$w$, the finite census estimate (obtained by applying the marked
estimate to the product over all $K$ slots) is
\[
 \|\Omega\|w^K\leq
   \exp\!\left(|\mathrm{Site}(k,p)|\,\mathrm{Treasury}_k+5K\right),
 \qquad
 \Omega=\int\prod_{j\in\mathrm{Site}(m,p)}U_j\,d\mu_{N,p}.
\]
Here $U_j$ is the $j$th translated/reflected census copy, with the chosen
motion applied, of the residual field $b-C.O(v)$. Thus $\Omega$ is exactly
the moment in the definition of the displayed root test.
For each slot, $x_k$ is the leading coefficient of its translated coefficient family
$\mathcal Z_k(-a)$ and $c_k$ its bill, as above, and the marched weight
has $w_k=w$. The rate clause at the opened scale $k$ is
$w_k\|\operatorname{alto}(H,k,d)\|(c_k+\|x_k\|_1)\leq\zeta_{d,k}$, and
$\zeta_{d,k}C_k\leq1$ by (R5). Hence
\[
 w_k\|\operatorname{alto}(H,k,d)\|C_k
       (c_k+\|x_k\|_1)\leq1,
 \qquad C_k=\operatorname{columns}(k)\geq1.
\]
The local hull estimate bounds the magnitude of its exponential entry by
$(1+C_k)\|x_k\|_1+c_k$: the unit coordinate costs at most $\|x_k\|_1$,
the masked core section at most $C_k\|x_k\|_1$ by the selected-family
price clause, and the masked regular part at most the bill $c_k$. Since
$(1+C_k)\|x_k\|_1+c_k\leq2C_k(c_k+\|x_k\|_1)$, the weighted entry price
is at most $2$. The forced profiles have total modulus at most $c_k$, so
the weighted total forced price is at most $1$. The action allowance is
$f=1$. Thus the marked estimate \eqref{ex-centered-sail} of
Subsection~\ref{ex-centered-construction} charges $x+y+2f=5$ per slot (the
more general choice $x=3,y=1$ gives the $6$-unit bound). The constant
$5K$ is exactly the sum of the per-slot tolls, and the extensive term is the cost
of all scale-$k$ sites. The aligned side-count identity is
\[
 |\mathrm{Site}(k,p)|=K(a_m/a_k)^4.
\]
Consequently, with
$g=(a_m/a_k)^4\mathrm{Treasury}_k+5$, we have
$\|\Omega\|w^K\leq e^{Kg}$. If $\Omega=0$ the desired inequality is
immediate; otherwise the same inequality still gives
$\|\Omega\|^{1/K}\leq e^g/w$. By definition the normalized census root
test is $\|\Omega\|^{1/K}$, hence
\[
 \operatorname{rootTest}_{m;N,p}(s\cdot(b-C.O(v)))
       \leq \frac{\operatorname{stakes}(m,k)}{w}.
\]
Finally, the definitions of escrow, flight, and the stake convert this
ratio exactly:
\[
 \frac{\operatorname{stakes}(m,k)}{w}
 =\operatorname{stakes}(m,k)
   \frac{4\|\operatorname{alto}(H,k,d)\|}{\zeta_k}
   \operatorname{escrow}(N,k,E)
 =\operatorname{cider}(d,m,k)E\operatorname{flight}_{-\gamma}(N,p).
\]
The bound holds for every $k$, $m$ and motion in the stated ranges,
because the record (R1)--(R7) and the Cream do not depend on them. No
subsequence or replacement history enters the construction.
\end{proof}

\subparagraph{Boundary of this producer.}
The proof above expands the terminal recentering, finite backward induction,
cost conversion, marched-weight calculation, and the final use of the
finite-copy estimate. It does not derive the selected history's Tone/Coupon
producer, its all-volume height ledger, or the selected-history binding of
the terminal Cream. Those exact suppliers are listed separately in the
evidence record; none follows merely from the existence of a portfolio.
The precise remaining certificate is the record (R1)--(R7); on one fixed
family and one fixed field of adjacent gauge doors, its same-history part
(R6)--(R7) reads as follows. For each degree $d$, terminal level
$N\geq\operatorname{ref}(d)$, and period $p$, it must provide a transfer sequence whose
terminal coefficient family is the chosen factory representative, whose played fields
realize the factory entries, which is admissible at every level from
$\operatorname{ref}(d)$ to $N$, and which responds across each intervening
door. The same transfer sequence must have Carol witnesses at all those doors and the
terminal forced-profile hook. For every $i\geq\operatorname{ref}(d)$ it
must also provide the paid local Coupon and its all-displacement Carol/Tale
comparison, with the uniform height, room, bank, Brine and monotone
occupancy prices used above. These are same-family, same-history
quantifiers: the transfer sequence is chosen for each $(d,N,p)$ and then retained
through all tests at those indices, whereas choosing a different history
or door field for each period or terminal volume is insufficient.

The nested field and adjacent comparisons are supplied by
Propositions \ref{prop:selected-terminal-voice},
\ref{prop:selected-reprise} and \ref{prop:selected-music}; the
zero-coordinate response is Proposition~\ref{prop:selected-zero-coordinate-coupon}.
Subsection~\ref{sec:selected-centered-service} identifies the height,
room, bank and occupancy prices in (R1)--(R7). The terminal Cream
attachment uses Proposition~\ref{prop:terminal-marked-voice}, including
its reserved-address, lawfulness and fee hypotheses. All these applications
retain the same original funded history as in
Proposition~\ref{prop:ex-original-history-attachment}; none permits a
new history to be chosen for each cutoff or period.
